\chapter{The example and the two controls}\label{ch:example}

\begin{chapterplan}
The paper's one example, Corollary~1 (the band of slots in which Proposition~2 applies), and two controls that
locate the rise in entry in the information in the price.
\end{chapterplan}

% ====================================================================
\section{The auction in the example}\label{sec:example-auction}

The example (Section~4.2 of the paper) shows that (H1)--(H3) can hold together. Class \(H\) has the one value \(h\), class \(L\) the one value \(\ell\), the reserve is \(p\), and the incumbent's value \(R\) is uniform
on \([0,r]\). On the path,
\begin{equation}\label{eq:example-order}
0<p<\ell<r<h.
\end{equation}
A larger \(r\) is a stronger incumbent, and class \(H\) always wins. The strengths are \(r_0<r_1<r_2\), the prior is
\(\tfrac12\), and the slot has \(1+u=\eu\) (Chapter~\ref{app:numbers}). The two classes pay differently only when
\(R>\ell\): the incumbent then pays \(\ell\) against class \(L\), while class \(H\) pays \(R\). So
\(\DT=\E[(R-\ell)^+]=(r-\ell)^2/(2r)\). The six objects, equation (8) of the paper, are
\begin{equation}\label{eq:example-payoffs}
\begin{aligned}
t_0&=p\Bigl(1-\frac pr\Bigr),& t_H&=\frac r2+\frac{p^2}{2r},& t_L&=\ell-\frac{\ell^2-p^2}{2r},\\
g_H&=h-\frac r2-\frac{p^2}{2r},& g_L&=\frac{\ell^2-p^2}{2r},& \DT(r)&=\frac{(r-\ell)^2}{2r}.
\end{aligned}
\end{equation}
Their derivatives (equation (9) of the paper) are
\(\DT'=\frac{r^2-\ell^2}{2r^2}>0\), \(g_H'=-\frac{r^2-p^2}{2r^2}<0\) and \(g_L'=-\frac{\ell^2-p^2}{2r^2}<0\). So
gross profit falls strictly at every belief, and the spread rises strictly. The gap \(\gamma=g_H-g_L\) satisfies
\(\DT+\gamma=h-\ell\). The threshold \(\tau=(c_H-g_L)/\gamma\) rises strictly on \(\ell\le r\le h\) (Appendix~A.6 of
the paper):
\begin{equation}\label{eq:example-tau}
\tau(r)=\frac{2rc_H-(\ell^2-p^2)}{2rh-r^2-\ell^2}.
\end{equation}

\section{Corollary 1: the band of slots}\label{sec:example-corollary}

Write \(u_{\rm lo}=\frac{2\tau(r_1)-1}{1-\tau(r_1)}\), \(u_{\rm hi}=\frac{2\tau(r_2)-1}{1-\tau(r_2)}\) and
\(u_k=\frac{\rho\,\DT(r_1)}{k}-2\).

\begin{paperresult}{Corollary 1 (Recruitment and ownership in the example, full statement)}
Let the classes be point masses at \(h\) and \(\ell\), the incumbent be uniform on \([0,r]\), and let
\((p,\rho,c_H,k)\) and \((r_0,r_1,r_2)\) take the example's values. Take a slot with \(0<\lambda<1\),
\(0<\beta<1/2\), \(u=\lambda/[(1-\lambda)\beta]\).
\begin{enumerate}[label=(\alph*),leftmargin=2em]
\item At prior \(\pi=\tfrac12\), the high-cost window, trading-cost window and exclusion condition hold together
  exactly on \(u_{\rm lo}<u<u_{\rm hi}\), equivalently \(\tau(r_1)<M(\pi)<\tau(r_2)\), with the endpoints of the band
  of Chapter~\ref{app:numbers}. This is an interval of presence probabilities at each fixed liquidity share. Under
  all four tie rules, entry is \(\rho\), \(\rho+(1-\rho)(b+\lambda/2)\), \(\rho\) at the three strengths, and
  high-class ownership is \(\rho/2\), \([\rho+(1-\rho)(b+\lambda)]/2\), \(\rho/2\).
\item For \(u<u_{\rm lo}\), entry at \(r_1\) is \(\rho\). At \(u=u_{\rm lo}\), orders are uniquely zero at \(r_0\)
  and uniquely full at \(r_1,r_2\). When an indifferent high-cost challenger prepares, every conclusion of
  Proposition~2 holds. When it may prepare with any probability, entry after a purchase at \(r_1\) ranges over
  \([\rho,1]\), and total entry there ranges over \([\rho,\rho+(1-\rho)(b+\lambda/2)]\). For
  \(u_{\rm hi}\le u<u_k\), a purchase at \(r_2\) induces certain preparation when an indifferent high-cost
  challenger prepares. For \(u\ge u_k\), the strict upper inequality of the trading-cost window at \(r_1\) fails.
\item If \(1+u=\eu\), all three hypotheses hold exactly when
  \[
  \frac{\tau(r_1)}{\tau(r_1)+[1-\tau(r_1)]\eu}<\pi<\frac{\tau(r_2)}{\tau(r_2)+[1-\tau(r_2)]\eu},
  \]
  with the endpoints of Chapter~\ref{app:numbers}. The high-cost and trading-cost windows alone hold exactly
  between the same lower endpoint and \(\tau(r_0)\).
\item At the example's slot (\(1+u=\eu\) and the example's \(\beta\)) and \(\pi=\tfrac12\), entry, high-class
  ownership and class \(H\)'s share of the acquisitions won by the challenger take, at the three strengths, the
  values of Chapter~\ref{app:numbers}.
\item For these point values and the strength range \([r_0,r_2]\), decisive advantage on the full uniform family
  and the affordability condition \(h-r_2\ge c_H\) of Proposition~3(c) hold together exactly on the cost window
  \[
  \frac{h^2-p^2}{2h}\le c_H\le h-r_2 .
  \]
  The example's cost lies in this window, with positive margins. For these class values, reserve and cost, every
  economy with a uniform incumbent \(U[0,r]\) in which some order recruits a high-cost challenger has
  \[
  r<h-c_H+\sqrt{(h-c_H)^2-p^2}<h,
  \]
  whatever the slot, prior, \(\rho\), \(k\) and research.
\end{enumerate}
\end{paperresult}

\paragraph{Proof.} At \(\pi=\tfrac12\), \(M=(1+u)/(2+u)\) and \(\kappa=\min\{m,1-M\}=1/(2+u)\). The left sides of (H1) and
(H2) do not contain the slot and hold at the example's values. The right side of (H1) is \(M>\tau(r_1)\), and
\[
\tfrac{1+u}{2+u}>\tau\iff 1+u>\tau(2+u)\iff u>\tfrac{2\tau-1}{1-\tau}.
\]
The same steps turn (H3) into \(u<u_{\rm hi}\). The right side of (H2) is \(u<u_k\), and \(u_{\rm hi}<u_k\), so it adds
nothing. For (c), \(M(\pi)>\tau\) holds exactly when the odds of \(\pi\) exceed \(\tau/[(1-\tau)(1+u)]\). For (e),
\(g_H(r)=\int_p^hF_r(y)\,dy\) falls strictly and equals \((h^2-p^2)/(2h)\) at \(r=h\); recruitment needs \(g_H>c_H\),
which is \(r<h-c_H+\sqrt{(h-c_H)^2-p^2}\), and this root lies below \(h\) exactly when \(c_H>(h^2-p^2)/(2h)\).
In words: at the prior \(\tfrac12\), the conditions of Proposition~2 contain the slot only through \(u\).

\paragraph{The entry ceiling.} The ceiling \(r_C\) solves \(\tau(r_C)=M(\tfrac12)\). On the band it is unique and lies
in \((r_1,r_2)\), just below \(r_2\). Left of \(r_C\) a purchase recruits; right of it nothing does
(Figure~\ref{fig:example-beliefs}, the paper's Figure~1). The band is narrow because the threshold rises slowly.
Differentiate \(g_L+\tau\gamma=c_H\) in \(r\):
\[
\tau'(r)=\frac{(1-\tau)(-g_L')+\tau(-g_H')}{\gamma}\in\Bigl(0,\frac{1}{2\gamma}\Bigr),
\]
because \(-g_H'=t_H'\) and \(-g_L'\) lie in \((0,\tfrac12)\). Class \(H\)'s payment rises by less than one half per unit
of strength, which is small next to the gap \(\gamma\).

\begin{figure}[htbp]
\centering
\includegraphics[width=0.75\linewidth]{beliefs}
\caption{Beliefs and the entry ceiling in the example, at the prior \(\tfrac12\). Panel (a): the preparation
threshold \(\tau(r)\) rises with strength from \(r_0\) to \(r_2\); the grey band holds every posterior between
\(m(\tfrac12)\) and \(M(\tfrac12)\); the threshold meets the top of the band at the entry ceiling \(r_C\). Panel (b):
the strengths from \(r_1\) to \(r_2\) on a finer vertical scale; the shaded band holds every top posterior that puts
\(r_C\) between \(r_1\) and \(r_2\). The printed values are those of Chapter~\ref{app:numbers}.}
\label{fig:example-beliefs}
\end{figure}

% ====================================================================
\section{The fixed-information control}\label{sec:example-control}

Proposition~2 changes the strength and the information in the price together. A control holds one fixed.

\begin{paperresult}{Proposition A.2 (Fixed information)}
Fix the joint law of the class \(j\) and a signal \(S\) that the challenger observes. Fix the class value laws
\(G_j\) and the cost law. In each economy \(i\), require that, conditional on \((j,S)\), the joint law of
\((\theta,R)\) is \(G_j\otimes F_i\). The cost is independent of the class, both values and \(S\). If the incumbent
strengthens from \(F_0\) to \(F_1\), entry is nonincreasing when an indifferent high-cost challenger prepares. Its
decrease is exactly
\[
\Ent_0-\Ent_1=\Pr\bigl\{B_{F_1}(\mu(S))<C\le B_{F_0}(\mu(S))\bigr\},\qquad \mu(S)=\Pr(j=H\mid S),
\]
and the decrease is strict exactly when this probability is positive.
\end{paperresult}

\paragraph{Proof.} The cost carries no news about the class, so the challenger prepares exactly when
\(C\le B_{F_i}(\mu(S))\). Proposition~1(a) gives \(B_1:=B_{F_1}(\mu(S))\le B_0:=B_{F_0}(\mu(S))\), and \((S,C)\) has
the same law in both economies. So
\(\Ent_0-\Ent_1=\E[\1\{C\le B_0\}-\1\{C\le B_1\}]=\Pr\{B_1<C\le B_0\}\ge0\). In words: at fixed information, a
stronger incumbent only lowers gross profit, so it can only lower entry.

\begin{watchout}
In the slot, fix full orders at every strength. The posteriors \((m,\pi,M)\) then do not depend on \(r\), but the
prices reveal them only if \(\DT>0\) in each economy. In the example \(\DT>0\) above \(\ell\).
\end{watchout}

\paragraph{The two panels.} Panel~B of the paper's outcome table imposes full orders: entry is
\(\rho+(1-\rho)(b+\lambda/2)\) at both \(r_0\) and \(r_1\). Panel~C hides the price: \(B_{r_1}(\pi)\le
B_{r_0}(\pi)<c_H\), so entry stays at \(\rho\). So the rise comes from the information that trading puts into the price.

% ====================================================================
\begin{chaptersummary}
In the example, gross profit falls and \(\DT=(r-\ell)^2/(2r)\) rises with strength. At the prior \(\tfrac12\),
(H1)--(H3) hold exactly when \(\tau(r_1)<M<\tau(r_2)\), a narrow band since \(\tau'<1/(2\gamma)\). The two controls
locate the rise of entry in the information that trading creates.
\end{chaptersummary}
